%\SweaveUTF8




\setkeys{Gin}{width=\textwidth}

\begin{Problem}[Population Characteristics]\label{problem:characteristics}
    We use the dataset \code{sepe.dat}, which contains data for a
    fictitious (sub-) population of firms:

\begin{Schunk}
\begin{Sinput}
> sepe <- read.table("http://stat.ethz.ch/Teaching/Datasets/WBL/sepe.dat",
                     header = TRUE)
\end{Sinput}
\end{Schunk}


    The dataset contains the number of full-time employees (\code{empl}),
    environmental protection expenditures (\code{expend}), and an indicator for
    industry class 31 (\code{code = 1}).

\begin{Subtask}
    %---------------------------------------------------------------------------
    \item Compute the total environmental protection expenditures, the mean
        expenditures per firm, the expenditures per employee in the population,
        and the mean of the firm-level expenditures per employee.

    \begin{Solution}
        We attach the data to the search space.

\begin{Schunk}
\begin{Sinput}
> attach(sepe)
\end{Sinput}
\end{Schunk}

        The requested quantities are:

        \begin{itemize}
            \item Total expenditures:

\begin{Schunk}
\begin{Sinput}
> sum(expend)
\end{Sinput}
\begin{Soutput}
[1] 68765270
\end{Soutput}
\end{Schunk}

            \item Mean expenditures per firm:

\begin{Schunk}
\begin{Sinput}
> mean(expend)
\end{Sinput}
\begin{Soutput}
[1] 85956.59
\end{Soutput}
\end{Schunk}

            \item Expenditures per employee in the population:

\begin{Schunk}
\begin{Sinput}
> sum(expend) / sum(empl)
\end{Sinput}
\begin{Soutput}
[1] 446.0035
\end{Soutput}
\end{Schunk}

            \item Mean firm-level expenditures per employee:

\begin{Schunk}
\begin{Sinput}
> mean(expend[empl > 0] / empl[empl > 0])
\end{Sinput}
\begin{Soutput}
[1] 1205.779
\end{Soutput}
\end{Schunk}
        \end{itemize}

        \begin{Remark}
            Three of the 800 firms have no recorded full-time employees, so
            their firm-level expenditures per employee are undefined. If these
            zeros are treated as missing values, the last mean can be computed
            from firms with positive recorded employment. This implicitly
            assumes that the omitted firms would have the same average ratio as
            the observed firms.
        \end{Remark}
    \end{Solution}

    %---------------------------------------------------------------------------
    \item Compute the population variance of the number of employees
        $\varu_{\mathrm{empl}}$, the population variance of expenditures
        $\varu_{\mathrm{expend}}$, and their correlation.

    \begin{Solution}
        Since the dataset contains the full population, use the population
        variance with denominator $N$:

        \begin{itemize}
            \item Variance of the number of employees:

\begin{Schunk}
\begin{Sinput}
> mean((empl - mean(empl))^2)
\end{Sinput}
\begin{Soutput}
[1] 36995.87
\end{Soutput}
\end{Schunk}

            \item Variance of expenditures:

\begin{Schunk}
\begin{Sinput}
> mean((expend - mean(expend))^2)
\end{Sinput}
\begin{Soutput}
[1] 4632073291
\end{Soutput}
\end{Schunk}

            \item Correlation between the number of employees and expenditures:

\begin{Schunk}
\begin{Sinput}
> cor(empl, expend)
\end{Sinput}
\begin{Soutput}
[1] 0.6597878
\end{Soutput}
\end{Schunk}
        \end{itemize}
    \end{Solution}

    %---------------------------------------------------------------------------
    \item Compute the proportion $\propu_{31}$ of firms in industry class 31.

    \begin{Solution}
        Since \code{code} is a binary indicator, the proportion is its mean:

\begin{Schunk}
\begin{Sinput}
> mean(code)
\end{Sinput}
\begin{Soutput}
[1] 0.1725
\end{Soutput}
\end{Schunk}
    \end{Solution}
\end{Subtask}
\end{Problem}
